\section{Математические основы AI-направленного прогнозирования термоупругих волн в геологических средах}

\subsection{Общие положения термоупругой волновой динамики}

Распространение термоупругих волн в геологических средах описывается системой связанных уравнений термоупругости, объединяющих уравнения движения и уравнения теплопереноса. В классической линейной модели Куна–Томпсона используется система:

\begin{equation}
\rho \frac{\partial^2 u_i}{\partial t^2} = \frac{\partial \sigma_{ij}}{\partial x_j},
\label{eq:motion}
\end{equation}

где $u_i$ — компоненты вектора смещения, $\rho$ — плотность среды, $\sigma_{ij}$ — компоненты тензора напряжений.

Термоупругая связь между деформациями и напряжениями задаётся:

\begin{equation}
\sigma_{ij} = \lambda \delta_{ij} \epsilon_{kk} + 2\mu \epsilon_{ij} - \gamma \delta_{ij} T,
\label{eq:stress}
\end{equation}

где $\lambda, \mu$ — параметры Ламе, $\epsilon_{ij}$ — тензор деформаций, $\gamma$ — коэффициент термоупругой связи, $T$ — изменение температуры.

Уравнение теплопереноса с учетом термомеханической генерации записывается как:

\begin{equation}
k \nabla^2 T - C \frac{\partial T}{\partial t} = \gamma T_0 \frac{\partial \epsilon_{kk}}{\partial t},
\label{eq:heat}
\end{equation}

где $k$ — коэффициент теплопроводности, $C$ — теплоёмкость, $T_0$ — средняя температура.

Данная система уравнений является базовой для численного моделирования и построения обучающих выборок для алгоритмов машинного обучения.

\subsection{Волновое уравнение в неоднородной геологической среде}

Геологические среды обладают выраженной неоднородностью: слоистостью, анизотропией, вариацией температуры и плотности. В случае пространственных вариаций параметров Ламе уравнение движения  принимает вид:

\begin{equation}
\rho(x) \frac{\partial^2 u_i}{\partial t^2} =
\frac{\partial}{\partial x_j}
\left[
\lambda(x) \delta_{ij} \frac{\partial u_k}{\partial x_k}
+ \mu(x) \left(
\frac{\partial u_i}{\partial x_j} + \frac{\partial u_j}{\partial x_i}
\right)
\right]
- \gamma(x) \frac{\partial T}{\partial x_i}.
\end{equation}

Анизотропия может быть учтена через тензор жесткости $C_{ijkl}$:

\begin{equation}
\sigma_{ij} = C_{ijkl} \epsilon_{kl} - \gamma_{ij} T.
\end{equation}

Моделирование таких сред требует высокопроизводительных численных методов: FDTD, спектральных методов или гибридных схем.

\subsection{AI-направленная инверсия и прогнозирование}

AI-направленный подход основан на интеграции моделей глубокого обучения с физически информированными системами PDE.  

Типовые методы включают:

\begin{itemize}
    \item физически-информированные нейронные сети (PINN);
    \item гибридные модели «симулятор + AI»;
    \item операторное обучение (Fourier Neural Operator, DeepONet);
    \item сверточные и трансформерные модели для прогноза термоупругих полей.
\end{itemize}

PINN обучается минимизируя функционал невязки дифференциальных уравнений:

\begin{equation}
\mathcal{L} = 
\left\|
\rho \frac{\partial^2 u_i}{\partial t^2} -
\frac{\partial \sigma_{ij}}{\partial x_j}
\right\|^2
+
\left\|
k \nabla^2 T - C \frac{\partial T}{\partial t}
- \gamma T_0 \frac{\partial \epsilon_{kk}}{\partial t}
\right\|^2
+
\mathcal{L}_\text{data},
\end{equation}

что обеспечивает физическую достоверность решения даже при ограниченном количестве измерений.

\subsection{Аппроксимация волнового оператора AI-моделью}

Операторное обучение позволяет аппроксимировать отображение:

\[
\mathcal{G} : 
\{ \lambda(x), \mu(x), \rho(x), T_0(x) \}
\rightarrow
\{ u(x,t), T(x,t) \},
\]

что эквивалентно построению быстрой ML-версии численного симулятора.  

Fourier Neural Operator формально реализует преобразование:

\begin{equation}
u^{(l+1)}(x) = \sigma\left( 
W u^{(l)}(x) + \mathcal{F}^{-1} \left[ R(k) \cdot \mathcal{F}(u^{(l)}) \right]
\right),
\end{equation}

где $\mathcal{F}$ — преобразование Фурье, $R(k)$ — обучаемая матрица в частотной области.  

Это позволяет прогнозировать волновые поля в десятки–сотни раз быстрее традиционных PDE-солверов.

\subsection{Оценка точности и метрические показатели}

Для оценки точности AI-предсказаний применяются:

\begin{itemize}
    \item $L_2$-норма ошибки:  
    \[
    E_{L_2} = \frac{\|u_{\text{AI}} - u_{\text{ref}}\|_2}{\|u_{\text{ref}}\|_2};
    \]
    \item структурное сходство (SSIM);
    \item энергия расхождения волн;
    \item временная ошибка фазового сдвига.
\end{itemize}

Подобные метрики применяются в задачах прогнозирования волнового поля в сейсмике.

\subsection{Заключение}

Математический аппарат термоупругих волн, интегрированный с современными AI-подходами, позволяет строить высокоточные модели прогноза поведения геологических систем. Совмещение PDE-моделирования и глубоких нейронных сетей открывает путь к быстрому моделированию и инверсии сложных подповерхностных структур.

\bibliographystyle{unsrt}

\begin{thebibliography}{10}

\bibitem{kuhn1966thermoelasticity}
Kuhn, W. (1966).
\textit{Thermoelasticity Theory}. Springer.

\bibitem{carcione2015wavefields}
Carcione, J. (2015).
\textit{Wave Fields in Real Media}. Elsevier.

\bibitem{virieux1986p}
Virieux, J. (1986).  
P– and S–wave FDTD modeling.  
\textit{Geophysics}.

\bibitem{sun2019thermoelasticFDTD}
Sun, Q., et al. (2019).  
Thermoelastic wave simulation using FDTD.  
\textit{Applied Thermal Engineering}.

\bibitem{raissi2019pinn}
Raissi, M., Perdikaris, P., Karniadakis, G. (2019).  
Physics-Informed Neural Networks.  
\textit{Journal of Computational Physics}.

\bibitem{li2020fno}
Li, Z., Kovachki, N., et al. (2020).  
Fourier Neural Operator.  
\textit{NeurIPS}.

\bibitem{bennett2023transformer}
Bennett, M. et al. (2023).  
Transformer models for wave propagation.  
\textit{Geophysical Journal International}.

\bibitem{wang2021deeplearningseismic}
Wang, S. (2021).  
Deep learning for seismic modeling.  
\textit{Reviews of Geophysics}.

\end{thebibliography}
